% Version fixed135 standalone complete insertion.
% Odd-degree 3-power towers: actual tame Galois quotients, distinct
% sigma vs sigma_2 coset actions, pro-2 group relator Nielsen
% elimination, rank-(3^k+2) quadratic Magnus forms, explicit
% nine-sheet Fox-Shapiro certificate, nested integral Hecke projectors.
%
% This is a marked finite-quotient theorem, NOT a globally canonical
% Fox gauge for arbitrary odd-index local field extensions.

\subsection{The $3$-power tame tower: marked Schreier--Fox
reductions and nested integral Hecke projectors}
\label{subsec:dyadic-f135-threepower-tower}

Fixed133 treated one index-three nonnormal Sylow preimage.
At index nine a new obstruction to naive copying appears:
the tame Frobenius has order six on the cosets.
Its $2$-primary part acts by inversion, while
the inverse of the whole Frobenius acts by
multiplication by $2^{-1}=5$ modulo nine.
The two operations coincide modulo three but
\emph{not} modulo nine.
Keeping these markings distinct gives an
actual tower of finite Schreier--Fox
certificates, compatible with the
nested odd-index transfer operators.

For $k\ge1$, put
\begin{equation}\label{eq:dyadic-f135-tower-fields}
n=3^k,\quad
F_n=\mathbf Q_2(2^{1/n}),\quad
L_n=F_n(\zeta_n),\quad
\Gamma=G_{\mathbf Q_2},\quad H_n=G_{F_n}.
\end{equation}
Choose compatible positive real algebraic
$n$th roots so that $F_3\subset F_9\subset\cdots$.
The polynomial $X^n-2$ is Eisenstein at two.
Moreover $\zeta_n$ is unramified of degree
\[
f_n=\operatorname{ord}_n(2)=2\cdot3^{k-1}
\]
over $\mathbf Q_2$.  The tame finite quotient
is therefore
\begin{equation}\label{eq:dyadic-f135-tame-quotient}
\operatorname{Gal}(L_n/\mathbf Q_2)
\simeq C_n\rtimes_{2} C_{f_n},\qquad
|\operatorname{Gal}(L_n/\mathbf Q_2)|=nf_n.
\end{equation}
The $\sigma,\tau$ marking of the
Roe--Turturean full-group source
\cite{RoeTurtureanGQ2} acts on the
$n$ right cosets $H_n\tau^i$ by
\begin{equation}\label{eq:dyadic-f135-coset-actions}
\boxed{\quad
i\xmapsto{\tau}i+1,\qquad
i\xmapsto{\sigma}2i,\qquad
i\xmapsto{\sigma^{-1}}\overline2\,i,
\qquad
i\xmapsto{\sigma_2}-i
\pmod n,
\quad}
\end{equation}
where $\overline2=2^{-1}\pmod n$ and
$\sigma_2=\sigma^{\omega_2}$.
Indeed the $2$-primary component of
the order-$f_n$ Frobenius is
$\sigma^{3^{k-1}}$, and
$2^{3^{k-1}}\equiv-1\pmod{3^k}$.
For $n=9$ the four actions in
\eqref{eq:dyadic-f135-coset-actions}
are, respectively,
$i+1$, $2i$, $5i$, and $-i$.
They must not be conflated.

\begin{theorem}[Uniform marked pro-$2$ Schreier presentation along the tower]
\label{thm:dyadic-f135-marked-schreier-tower}
For every $n=3^k$,
the maximal pro-$2$ Galois quotient
$P_n=G_{F_n}(2)$ has the following
explicit marked pro-$2$ presentation.
Start with the right-coset
representatives $t_i=\tau^i$.
The finite-index Schreier scan of
the genuine four-generator
Roe--Turturean marked group
has $3n+1$ generators before
the tree/relator changes:
\[
s_i=\tau^i\sigma
  \tau^{-(2i\bmod n)},\quad
t=\tau^n,\quad
w_{ij}=\tau^ix_j\tau^{-i}
\quad (0\le i<n,\ j=0,1).
\]
The $n$ indexed tame relations,
obtained by scanning
$\sigma^{-1}\tau\sigma\tau^{-2}$,
have an integral unit Fox minor.
On the maximal pro-$2$ quotient
they eliminate
$t$ and $n-1$ of the $s_i$,
leaving one generator $s$.

Write $a_i=w_{i0}$, $v_i=w_{i1}$
with indices in $\mathbf Z/n$.
After tame elimination put
\begin{equation}\label{eq:dyadic-f135-indexed-wild-packets}
\boxed{
\begin{aligned}
u_{j,i}&=
\left(\prod_{r=0}^{n-1}
 w_{i+r,j}\right)^{1/n},
&
d_i&=u_{0,i}a_i^{-1},
\\
z_i&=s^{-1}a_{-i}s,
&
c_i&=[d_i,z_i],
\\
g&=s^2,
&
D_i&=g^{-1}d_i g,
\\
h_i&=[D_i,d_i],
&
H_i&=(g^{-1}a_i g)a_iD_i d_i^3h_i .
\end{aligned}}
\end{equation}
Products are taken in the displayed
cyclic increasing order; the
continuous $1/n$ exponent is
well-defined in pro-$2$ groups.
The \emph{actual} $n$ indexed wild
Schreier relators are
\begin{equation}\label{eq:dyadic-f135-indexed-wild-relators}
\boxed{\qquad
R_i=H_i u_{1,i}^{-1}
\,s^{-1}v_{\overline2 i}s\,c_i,
\qquad i\in\mathbf Z/n.
\qquad}
\end{equation}
The two different indices $-i$ and
$\overline2i$ reflect the separate
$\sigma_2$ and $\sigma^{-1}$
actions of \eqref{eq:dyadic-f135-coset-actions}.
Consequently
\begin{equation}\label{eq:dyadic-f135-tower-n-two-n-relators}
\boxed{\quad
P_n\cong
\left\langle
s,a_0,\ldots,a_{n-1},
v_0,\ldots,v_{n-1}
\ \middle|\
R_0,\ldots,R_{n-1}
\right\rangle_{\mathrm{pro}-2}.
\quad}
\end{equation}
This statement is at the \emph{group
relation} level, not a specialization
of the coefficient exchange Laplacian
of fixed132.
\end{theorem}

\begin{proof}
The tame normal extension
\eqref{eq:dyadic-f135-tame-quotient}
acts transitively on the $n$ roots.
The right-coset Schreier formula is
$t_i g=h_{i,g}t_{j(i,g)}$.
The $\tau$ spanning tree has $n-1$
edges, and the remaining generating
edges are the $n$ displayed $s_i$,
one $t=\tau^n$, and the $2n$ wild loops.
The tame relation reduces the
subgroup generated by the $s_i,t$
to the pro-$2$ quotient of
\[
\langle \sigma,t\mid
\sigma^{-1}t\sigma=t^2\rangle .
\]
Because a nontrivial element of
a finite $2$-group cannot be
conjugate to its square,
$t=1$.  The other $s_i$ then
have the common image $s$.
The primitive tame Fox minor
is also proved explicitly in
Theorem~\ref{thm:dyadic-f135-n-sheet-fox-matrix}.

For the wild words, before killing $t$
one has the exact indexed identity
\[
\tau^i(x_j\tau)^n\tau^{-i}
=\left(\prod_{r=0}^{n-1}
\widetilde w_{i+r,j}\right)\tau^n,
\quad
\widetilde w_{h,j}=
\tau^h x_j\tau^{-h}.
\]
The extended indices differ from
their residues modulo $n$ by
conjugation with $t$.
After $t=1$ in the pro-$2$ quotient,
the image of
$(x_j\tau)^{\omega_2}$
is the unique pro-$2$ $n$th
root of this cyclic product,
giving $u_{j,i}$.

The $2$-primary tame component
$\sigma_2$ acts by $-i$ on
the cosets, whereas the
ordinary inverse Frobenius
in $x_1^\sigma=\sigma^{-1}x_1\sigma$
acts by $\overline2i$.
In the pro-$2$ quotient
the images of $\sigma_2$ and
$\sigma$ in the stabilizer
are the same $s$,
but their \emph{coset}
effects remain distinct.
Hence $x_0^{\sigma_2}$
rewrites as $s^{-1}a_{-i}s$,
and $x_1^\sigma$ as
$s^{-1}v_{\overline2i}s$.
The remaining marked auxiliary
words involve the displayed
products, conjugations and
commutators and give exactly
\eqref{eq:dyadic-f135-indexed-wild-packets}
and
\eqref{eq:dyadic-f135-indexed-wild-relators}.

The marked pro-$2$ normal-closure
condition persists: every finite
$2$-group quotient of $H_n$
induces a quotient of $\Gamma$
into the corresponding finite
wreath group with $n$ cosets;
the images of the wild marked
generators lie in its normal
$2$-group base.
Conversely the indexed
Schreier equations supply the
corresponding admissible marked
finite wreath representations.
By the finite-quotient
universality of the
Roe--Turturean presentation,
the two pro-$2$ finite-target
functors agree.
Taking their inverse limits
proves the claimed group
presentation, with no use
of the false bare
full-group two-relator
presentation.
\end{proof}

\begin{theorem}[Uniform primitive wild Nielsen reduction and quadratic form]
\label{thm:dyadic-f135-uniform-nielsen-demushkin}
The augmented Fox row of the
$n$ genuine pro-$2$ wild relations
of Theorem~\ref{thm:dyadic-f135-marked-schreier-tower}
is
\begin{equation}\label{eq:dyadic-f135-nwild-augmented-fox}
\boxed{
\begin{aligned}
\epsilon(\partial_sR_i)&=0,\\
\epsilon(\partial_{a_j}R_i)&=\frac4n-2\delta_{ij},\\
\epsilon(\partial_{v_j}R_i)
&=\delta_{j,\overline2 i}-\frac1n.
\end{aligned}}
\end{equation}
Deleting the $i=0$ relation row
and $v_0$ generator column gives
the exact wild $(n-1)\times(n-1)$
minor
\begin{equation}\label{eq:dyadic-f135-wild-n-minus1-minor}
\boxed{\qquad
\det\left(
\delta_{j,\overline2 i}-\frac1n
\right)_{1\le i,j<n}
=\pm\frac1n
\in\mathbf Z_2^\times.
\qquad}
\end{equation}
Therefore the $n-1$ relators
$R_1,\ldots,R_{n-1}$
may replace the generators
$v_1,\ldots,v_{n-1}$
in a genuine free pro-$2$
Nielsen automorphism.
Killing these primitive
relation letters leaves
the $n+2$ generators
\[
s,a_0,\ldots,a_{n-1},v_0
\]
and the single pro-$2$ relator
\begin{equation}\label{eq:dyadic-f135-uniform-one-relator}
\boxed{
\mathscr R_n:=
\operatorname{ev}_{R_1=\cdots=R_{n-1}=1}
\left(\alpha_n^{-1}
\left(\prod_{i=0}^{n-1}R_i\right)\right).
}
\end{equation}
Here $\alpha_n$ is the chosen
free pro-$2$ basis automorphism;
its inverse is obtained
compatibly in finite
pro-$2$ quotient stages.
The complete augmented
minimal Fox row is
\begin{equation}\label{eq:dyadic-f135-uniform-minimal-first-row}
\boxed{\qquad
\epsilon(\partial_s\mathscr R_n)=0,
\quad
\epsilon(\partial_{a_j}\mathscr R_n)=2,
\quad
\epsilon(\partial_{v_0}\mathscr R_n)=0.
\qquad}
\end{equation}

Write $S,A_0,\ldots,A_{n-1},V$
for the mod-two Magnus variables
of this minimal generating set.
The quadratic relator symbol is
\begin{equation}\label{eq:dyadic-f135-uniform-magnus-quadratic}
\boxed{\quad
\sigma_2(\mathscr R_n)
=\sum_{i=0}^{n-1}A_i^2+SV+VS.
\quad}
\end{equation}
The cup pairing therefore has
nonalternating rank $n+2$,
with matrix
\begin{equation}\label{eq:dyadic-f135-tower-cup-matrix}
\boxed{\qquad
B_n=
\begin{pmatrix}
0&0&1\\
0&I_n&0\\
1&0&0
\end{pmatrix},
\qquad
\det_{\mathbf F_2} B_n=1,
\qquad}
\end{equation}
in the ordered basis
$(s,a_0,\ldots,a_{n-1},v_0)$.
In particular the first Fox
ideal is $(2)$, and
\begin{equation}\label{eq:dyadic-f135-uniform-abelianization}
P_n^{\mathrm{ab}}
\simeq\mathbf Z_2^{\,n+1}
\oplus\mathbf Z/2.
\end{equation}
Thus the index-nine case
has eleven minimal generators
and the exact quadratic symbol
\[
\sum_{i=0}^{8}A_i^2+SV+VS.
\]
The group-level Fox complex is
the genuine length-two
Demu\v{s}kin resolution, not
an augmentation-only model.
\end{theorem}

\begin{proof}
In abelianization,
$[u_{j,i}]=
\frac1n\sum_h[w_{h,j}]$.
Exactly as in fixed133,
\[
[H_i]=2[a_i]+4([u_{0,i}]-[a_i])
=\frac4n\sum_h[a_h]-2[a_i],
\]
and the other factors give
$[v_{\overline2i}]-\frac1n\sum_h[v_h]$.
This proves the first row.

The permutation
$i\mapsto\overline2 i$
fixes zero and permutes
the nonzero residues.
Hence on indices $1,\dots,n-1$
the wild block equals
$P-\frac1nJ_{n-1}$
for a permutation matrix $P$.
Its determinant is
\[
\det(P)\det\left(
I-\frac1nJ_{n-1}\right)
=\det(P)\left(1-\frac{n-1}{n}\right)
=\pm\frac1n.
\]
As $n$ is odd, this is
a unit of $\mathbf Z_2$.
The pro-$2$ Burnside
basis theorem and the
Hopfian property of finite-rank
free pro-$2$ groups make the
specified substitution an
actual automorphism, just as
in fixed133.  The product
of the $n$ rows has
exact integral exponent
sum $2$ in each $a_i$,
zero in all $v_j$,
and zero in $s$.
This proves
\eqref{eq:dyadic-f135-uniform-one-relator}
and
\eqref{eq:dyadic-f135-uniform-minimal-first-row}.

For the degree-two symbol,
it is enough to identify
all $v_i$ at the linear
Magnus level with one
variable $V$.
Indeed the eliminated
wild relators equate
their linear classes;
the product of all $R_i$
has vanishing \emph{mod-two
linear} Magnus term,
so higher-order corrections
to these substitutions
cannot alter its quadratic
part.

After making the uniform
substitution $v_i=v$ as
group words, the cyclic
product defining $u_{1,i}$
is $v^n$, so its pro-$2$
$n$th root is \emph{exactly}
$v$, regardless of
$\binom{1/n}{2}\bmod2$.
The two factors
$u_{1,i}^{-1}
s^{-1}v_{\overline2i}s$
therefore reduce to
the commutator
$[v,s]$, whose degree-two
symbol is $VS+SV$.
The $g=s^2$ conjugations
and fourth powers of the
$d_i$ contribute no
degree-two terms, so
$H_i$ contributes
$A_i^2$ modulo degree three.

Finally let
$A_\Sigma=\sum_i A_i$.
The linear Magnus term
of $d_i$ is
$A_\Sigma+A_i$;
the linear term
of $z_i=s^{-1}a_{-i}s$
is $A_{-i}$.
Thus the total degree-two
commutator contribution
from all $c_i$ is
\[
\sum_i[A_\Sigma+A_i,A_{-i}]
=[A_\Sigma,A_\Sigma]
+\sum_i[A_i,A_{-i}]=0
\]
in characteristic two:
the terms $i$ and $-i$
cancel in pairs, and
$i=0$ gives a trivial
self-commutator.
Each folded $R_i$ starts
in Magnus degree at least
two.  Multiplying the
$n$ relators therefore
adds their quadratic parts.
As $n$ is odd,
$n[V,s]=[V,s]$,
proving
\eqref{eq:dyadic-f135-uniform-magnus-quadratic}.
The cup pairing is the
ordered quadratic Fox
symbol, so the displayed
matrix is immediate.
Its determinant is one,
and the integral first
Fox row has ideal $(2)$,
giving the abelianization.

Since $F_n/\mathbf Q_2$
is a finite dyadic field
of degree $n$ containing
$\mu_2$ but not $\mu_4$,
the maximal pro-$2$ local
Galois group is a
Demu\v{s}kin group of
rank $n+2$.
The resulting genuine
minimal one-relator
Fox complex is the
exact completed
projective resolution.
No extra Fox relation
syzygy is assumed.
\end{proof}

\begin{theorem}[A uniform $n$-sheet integral Fox--Shapiro certificate]
\label{thm:dyadic-f135-n-sheet-fox-matrix}
Let $n=3^k$ and let
$W_n=\mathbf Z_2[\Gamma/H_n]$
be the $n$-dimensional
coset permutation module.
The matrices $\Sigma,\Upsilon$
of the tame generators
are determined by
\begin{equation}\label{eq:dyadic-f135-coset-permutation-matrices}
\boxed{
\Sigma_{i,j}=\delta_{j,2i\bmod n},
\qquad
\Upsilon_{i,j}=\delta_{j,i+1\bmod n}.
}
\end{equation}
Put $I=I_n$,
$J=\mathbf1_n\mathbf1_n^t$,
$P_n=J/n$,
and let $e_{n-1}$ be
the final coordinate vector.
The full \emph{rationally
written but $2$-integral}
Roe--Turturean coefficient
differentials on
$W_n$ are
\begin{equation}\label{eq:dyadic-f135-full-n-sheet-fox-differentials}
\boxed{
\begin{aligned}
d^0(b)&=((\Sigma-I)b,
(\Upsilon-I)b,0,0),\\
d^1(a,b,c,d)
&=\bigl(
\Sigma^{-1}(\Upsilon-I)a+
(\Sigma^{-1}-I-\Upsilon)b,\\
&\hspace{1em}
3P_n b+(4P_n-2I)c+
(\Sigma^{-1}-P_n)d
\bigr).
\end{aligned}}
\end{equation}
The $n-1$ elementary
$\tau$-spanning-tree edges
are a strict contractible
vertex--edge block.
Remove their columns,
and multiply each
of the $n$ wild
relation rows by the
odd unit $n$.
The resulting
complete integer
$(2n)\times(3n+1)$
augmented Fox matrix,
with column blocks
$(s_0,\ldots,s_{n-1};
t;a_0,\ldots,a_{n-1};
v_0,\ldots,v_{n-1})$,
is
\begin{equation}\label{eq:dyadic-f135-n-sheet-augmented-matrix}
\boxed{
\mathsf J^{(n)}=
\begin{pmatrix}
\Sigma^{-1}(\Upsilon-I)
&-e_{n-1}&0&0\\[2pt]
0&3\mathbf1_n&
4J-2nI&
n\Sigma^{-1}-J
\end{pmatrix}.
}
\end{equation}
The tame $n\times n$
minor formed by columns
$s_1,\ldots,s_{n-1},t$
has determinant $\pm1$.
After this tame elimination,
the wild $n\times n$
block against all $v_i$
has rank $n-1$ modulo two,
and the $(n-1)\times(n-1)$
minor omitting row and column
zero has determinant
$\pm n^{n-2}$ in the
integer matrix.
The surviving
integral face row has
precisely one factor two.
Consequently
\begin{equation}\label{eq:dyadic-f135-uniform-smith-invariants}
\boxed{\quad
\operatorname{SNF}_{\mathbf Z_2}
(\mathsf J^{(n)})
=(1^{\,2n-1},2).
\quad}
\end{equation}
In particular for $n=9$,
$\mathsf J^{(9)}$ is
$18\times28$,
has mod-two rank $17$,
and the ordered
eighteen-column minor
\begin{equation}\label{eq:dyadic-f135-nine-exact-minor}
\boxed{
\begin{gathered}
\mathcal I_{18}
=(1,2,3,4,5,6,7,8,10,11,
20,21,22,23,24,25,26,27),\\
\det\bigl(\mathsf J^{(9)}
_{[1:18],\mathcal I_{18}}\bigr)
=-18\cdot9^7=-86{,}093{,}442 .
\end{gathered}}
\end{equation}
There is also a
seventeen-column minor
with column list
$(1,2,3,4,5,6,7,8,10,
20,21,22,23,24,25,26,27)$
on rows $1,\dots,17$
and determinant $-9^7$.
These two literal
integer determinants
certify the complete
local Smith valuation
without floating
point approximation.
\end{theorem}

\begin{proof}
Use the coinduced
finite coset module
and the left
crossed-derivation
Fox convention.
The tame relator
gives the first
component of
\eqref{eq:dyadic-f135-full-n-sheet-fox-differentials}.
The permutation
$\Upsilon$ has
odd order $n$,
so the $\omega_2$
projection of an
affine lift
$(\Upsilon,b+c)$
is the invariant
translation
$(I,P_n(b+c))$:
its $n$th power
is the sum over
the cyclic orbit,
and division by
$n$ is integral.
The $2$-part
$\sigma_2$ of
the affine tame
Frobenius has
permutation part
$\Sigma^{3^{k-1}}$,
whose square is
identity.
Its own square
has trivial
permutation part,
so it conjugates
pure translations
trivially.
All commutators
of the auxiliary
wild pure
translations vanish.
Consequently the
marked wild relator
reduces to
$2c+4(P_n(b+c)-c)
-P_n(b+d)+\Sigma^{-1}d$,
which simplifies
to the second
component of
\eqref{eq:dyadic-f135-full-n-sheet-fox-differentials}.

The tree coboundary
$(\Upsilon-I)v$
has $n-1$ primitive
edge differences.
Contract them,
retaining the last
$\tau$ edge.
The tame remaining
column is
$(\Sigma^{-1}-I-\Upsilon)
e_{n-1}=-e_{n-1}$,
and the wild
remaining column
after row
multiplication by
$n$ is
$3J e_{n-1}
=3\mathbf1_n$.
This gives
\eqref{eq:dyadic-f135-n-sheet-augmented-matrix}
exactly.
The upper tame
block is the
Fox boundary of
the $n$ tame
Schreier relators.
It is a cycle
difference matrix
with one
primitive $t$ pivot;
deleting any
one of the
$s_i$ columns
and retaining
$t$ gives a
determinant $\pm1$.

The lower wild
$v$ block is
$n\Sigma^{-1}-J$.
The permutation
$\Sigma^{-1}$
fixes index zero
and permutes the
other $n-1$
indices.
Its lower-right
$(n-1)\times(n-1)$
minor is therefore
\[
\det\bigl(
nP-J_{n-1}\bigr)
=\det(P)n^{n-1}
\left(1-\frac{n-1}{n}\right)
=\pm n^{n-2},
\]
a $2$-adic unit.
Hence there are
$n$ tame plus
$n-1$ wild
primitive Fox
pivots.
After their
elimination the
surviving face
is the product
$\prod_iR_i$
from
Theorem~\ref{thm:dyadic-f135-uniform-nielsen-demushkin}.
Its exact
augmented first
Fox row is two
in each of the
$n$ $a_i$ columns
and zero elsewhere,
so its last
Smith valuation
is exactly one.
This proves the
uniform Smith
formula.
For $n=9$,
substituting the
nine explicit
permutation
matrices and
taking ordinary
integer determinants
gives the two
declared values
in
\eqref{eq:dyadic-f135-nine-exact-minor}.
They also appear
as exact
regression assertions
in the companion
Sage source.
\end{proof}


\begin{theorem}[Uniform cyclotomic Fox row and primitive top-face trace]
\label{thm:dyadic-f135-uniform-fox-tate-trace}
In the actual $(1+2n)$-generator,
$n$-relation pro-$2$ presentation
\eqref{eq:dyadic-f135-tower-n-two-n-relators},
let $\theta$ denote the local
cyclotomic orientation.  Its marked
values are
\[
\theta(s)=1,\qquad
\theta(a_j)=-1,\qquad
\theta(v_j)=-1/3.
\]
Put
\[
c=-1/3,\qquad
\lambda_n=\frac{c-1}{c^n-1}
=\frac{4\cdot3^{n-1}}{3^n+1}
\in\mathbf Z_2^\times.
\]
For a $\theta$-crossed derivation
$D$, write $A_j=D(a_j)$ and
$V_j=D(v_j)$.
The complete twisted Fox row of
the $i$th indexed wild relator is
\begin{equation}\label{eq:dyadic-f135-uniform-twisted-fox}
\boxed{
\begin{aligned}
D(R_i)={}&
2\left(\sum_{r=0}^{n-1}
(-1)^rA_{i+r}-A_i\right)\\
&+3\lambda_n
\sum_{r=0}^{n-1}c^rV_{i+r}
-3V_{\overline2 i}.
\end{aligned}}
\end{equation}
In particular, the sum of these
$n$ twisted Fox rows vanishes
\emph{integrally, columnwise}.
The $(n-1)$-square wild $v$
minor is a unit modulo two, and
so the full twisted top-image is
exactly the kernel of the
primitive face sum
\begin{equation}\label{eq:dyadic-f135-uniform-face-trace}
\boxed{
\operatorname{Tr}_{F_n}:
\mathbf Z_2^n\longrightarrow
\mathbf Z_2,\qquad
(q_i)\longmapsto\sum_iq_i.
}
\end{equation}
For trivial coefficients the
same trace has image
$2\mathbf Z_2$ on top coboundaries.
Thus
\begin{equation}\label{eq:dyadic-f135-uniform-tate-ledgers}
\boxed{
\begin{array}{c|ccc}
 &H^0&H^1&H^2\\\hline
\mathbf Z_2&
\mathbf Z_2&
\mathbf Z_2^{n+1}&\mathbf F_2\\
\mathbf Z_2(1)&
0&\mathbf Z_2^{n+1}
\oplus\mathbf F_2&\mathbf Z_2 .
\end{array}}
\end{equation}
The Wu/Bockstein class in the
minimal marked basis is
\[
w_{F_n}=\sum_{i=0}^{n-1}
\chi_{a_i},\qquad
\chi^2=\chi\smile w_{F_n}.
\]
For $F_{3^j}\subset F_{3^k}$,
$j\le k$, the compatible
primitive cyclotomic top-face
traces satisfy
\begin{equation}\label{eq:dyadic-f135-top-tate-transitivity}
\boxed{
\operatorname{Tr}_{F_{3^k}}
\operatorname{res}_{F_{3^k}/F_{3^j}}
=3^{k-j}\operatorname{Tr}_{F_{3^j}}.
}
\end{equation}
The analogous trivial-coefficient
formula holds modulo two.
This is the exact odd-degree
local Tate invariant multiplier.
\end{theorem}

\begin{proof}
The orientation of a cyclic
$a$ packet is $-1$ and its
pro-$2$ $1/n$ root has the
same orientation; its twisted
Fox derivative factor is one.
A cyclic $v$ packet has
orientation $c^n$, and the
$1/n$ root has orientation $c$,
giving the scalar factor
$\lambda_n$.
For odd $n$,
$v_2(3^n+1)=2$;
hence $\lambda_n$ is integral
and is a $2$-adic unit.

Put $D_i=D(d_i)$.
By the inverse and product
Fox rules,
\[
D_i=\sum_r(-1)^rA_{i+r}-A_i.
\]
The other marked factors have
orientations
$\theta(d_i)=1$,
$\theta(z_i)=-1$ and
$\theta(u_{1,i})=c$.
Consequently
\[
D([d_i,z_i])=-2D_i,
\quad
D(H_i)=-4D(s)+4D_i.
\]
The inverse $u_{1,i}^{-1}$
contributes $3D(u_{1,i})$;
the conjugate
$s^{-1}v_{\overline2 i}s$
contributes
$-3V_{\overline2 i}+4D(s)$.
The two $D(s)$ terms cancel.
Adding the last commutator
gives the displayed row.

The alternating sum
$\sum_{r=0}^{n-1}(-1)^r$
equals one, so
$\sum_iD_i=0$.
Moreover
$\lambda_n\sum_r c^r=1$,
and $i\mapsto\overline2 i$
permutes $\mathbf Z/n$.
Summing the $n$ rows
therefore gives zero
on every generator.
Modulo two the orientation
is trivial, hence the
$v$-minor has the same
unit reduction as
\eqref{eq:dyadic-f135-wild-n-minus1-minor}.
The top twisted image is
the sum-zero hyperplane.
For trivial coefficients,
the unit minor gives the
same hyperplane while the
sum of the $a_j$ columns
is exactly two, giving the
preimage of $2\mathbf Z_2$.

For the twist the first
differential is
\[
1\longmapsto
(0,-2,\ldots,-2,
-4/3,\ldots,-4/3)^t,
\]
with one Smith factor two.
The twisted second
differential has $n-1$
unit factors and one zero
face direction.  For the
trivial complex, the
second differential has
$n-1$ units and one
factor two while $d^0=0$.
This proves the integral
cohomology ledgers.
The mod-four orientation
reduces to the three-power
sum of the $\chi_{a_i}$;
the nonalternating quadratic
symbol of
\eqref{eq:dyadic-f135-uniform-magnus-quadratic}
identifies it with the
Wu characteristic class.

Finally compatible Schreier
covers pull back a marked
top face to $3^{k-j}$
copies.  The primitive
relation-face sum therefore
multiplies by $3^{k-j}$.
The strict coinduced
idempotents of
Theorem~\ref{thm:dyadic-f135-nine-hecke-tower}
identify that source
operation with derived
restriction.  Since the
index is odd it agrees
with the local Tate
fundamental-class
normalization.
\end{proof}

\begin{theorem}[Strict nested odd-index projectors: ranks $1,2,6$]
\label{thm:dyadic-f135-nine-hecke-tower}
On the nine-sheet
module $W_9$, define
\begin{equation}\label{eq:dyadic-f135-nine-projection-data}
\begin{aligned}
(P_1)_{ij}&=\frac19,\\
(P_3)_{ij}&=
\begin{cases}
1/3,&i\equiv j\pmod3,\\
0,&i\not\equiv j\pmod3.
\end{cases}
\end{aligned}
\qquad (0\le i,j<9).
\end{equation}
Put
\begin{equation}\label{eq:dyadic-f135-nine-three-projectors}
\boxed{\quad
E_0=P_1,\qquad
E_1=P_3-P_1,\qquad
E_2=I_9-P_3.
\quad}
\end{equation}
These are integral
mutually orthogonal
idempotents over
$\mathbf Z_2$:
\begin{equation}\label{eq:dyadic-f135-nine-orthogonal}
\boxed{
E_iE_j=\delta_{ij}E_i,\quad
E_0+E_1+E_2=I_9,\quad
(\operatorname{rank}E_0,
 \operatorname{rank}E_1,
 \operatorname{rank}E_2)
=(1,2,6).
}
\end{equation}
They commute with
the full tame
coset permutation
matrices $\Sigma,\Upsilon$,
hence with
the entire
marked $\Gamma$-action
on the nine cosets.
For every finite
projective complete
$\Gamma$ coefficient
module $B$,
they give a
\emph{strict integral}
decomposition in
every degree of
the full nine-sheet
coinduced
Fox--Shapiro complex:
\begin{equation}\label{eq:dyadic-f135-nine-three-chain-summands}
\boxed{
\begin{aligned}
C^\bullet_{\mathrm{cover},9}
(H_9,B)
\cong {}&
\mathcal J_\Gamma^\bullet(B)\\
&\oplus\mathcal J_\Gamma^\bullet
(B\otimes T_3)\\
&\oplus\mathcal J_\Gamma^\bullet
(B\otimes W_6).
\end{aligned}}
\end{equation}
Here $T_3$ is the
rank-two augmentation
lattice of the
three-coset $S_3$
quotient of fixed134,
while $W_6$ is
the rank-six kernel
of the transfer
$W_9\to W_3$.
The three strict
chain idempotents are
$E_i$ in degree
zero, $E_i^{\oplus4}$
in degree one,
and $E_i^{\oplus2}$
in degree two,
acting coefficientwise
on $B$.

Under the
Fox--Shapiro comparison,
$P_3$ is the normalized
Hecke idempotent
\[
\frac13\operatorname{res}_{H_9}^{H_3}
\operatorname{cor}_{H_9}^{H_3}
\quad\text{on }R\Gamma(H_9,B),
\]
and $P_1$ is
\[
\frac19\operatorname{res}_{H_9}^{\Gamma}
\operatorname{cor}_{H_9}^{\Gamma}.
\]
They satisfy
\begin{equation}\label{eq:dyadic-f135-nested-hecke-identity}
\boxed{\quad
P_3P_1=P_1P_3=P_1,\qquad
(P_3-P_1)P_1=0.
\quad}
\end{equation}
Thus two successive
odd-index descents
agree with direct
index-nine descent
\emph{as strict
cochain projectors},
not merely
on cohomology.

For every residual
coefficient sector
on which $B|_{H_9}$
has pro-$2$ image,
Theorems~\ref{thm:dyadic-f135-marked-schreier-tower}
and~\ref{thm:dyadic-f135-uniform-nielsen-demushkin}
make the nine-sheet
cover a genuine
perfect $H_9$
Fox resolution.
Therefore the
three summands in
\eqref{eq:dyadic-f135-nine-three-chain-summands}
represent their
intrinsic derived
Galois cochain
objects.  In
particular the
constant summand
gives
\begin{equation}\label{eq:dyadic-f135-nine-full-marked-sector-completeness}
\boxed{\quad
\mathcal J_\Gamma^\bullet(B)
\simeq R\Gamma(\Gamma,B)
\quad\text{on this
index-nine residual
coefficient sector}.
\quad}
\end{equation}
This extends
fixed134's
index-three
coefficient completeness
to the genuinely
larger tame quotient
$C_9\rtimes C_6$.
It is not a
universal projective
two-relator resolution
of every absolute
Galois coefficient
sector.
\end{theorem}

\begin{proof}
The partition
of the nine
cosets by residues
modulo three is
preserved by
$i\mapsto i+1$
and $i\mapsto2i$.
The blockwise
coset average
$P_3$ is therefore
$\Gamma$-equivariant.
The total
average $P_1$
is plainly
$\Gamma$-equivariant.
Each is an
idempotent,
and
$P_3P_1=P_1$.
The orthogonal
idempotents
$E_0,E_1,E_2$
and their ranks
follow by direct
matrix multiplication.
As only the
odd scalars
$1/3$ and
$1/9$ are used,
this is integral
over $\mathbf Z_2$.

The transfer
$W_9\to W_3$
sums entries
over the three
cosets with the
same residue
modulo three.
Its section
sends one
three-coset basis
vector to the
sum of its
three lifts,
and the composite
is multiplication
by $3$.
Its normalized
idempotent is
exactly $P_3$.
The further
projection from
$W_3$ to the
constant line
is $P_1$ on
$W_9$.
The quotient
of the
three-coset
module by its
constant line
is the
rank-two natural
$S_3$ lattice
of fixed134.
The kernel of
the first
fiber sum has
rank six,
giving the
three summands
in
\eqref{eq:dyadic-f135-nine-three-chain-summands}.

For arbitrary
$B$ the
coinduced module
is, after the
standard transport
trivialization,
$B\otimes W_9$.
All these
idempotents are
$\Gamma$-equivariant.
Completed Fox
cochains are
functorial in
the coefficient
module, so their
degreewise
extensions commute
\emph{strictly}
with every
marked Fox
differential,
including the
profinite
$\omega_2$
operations.
This proves the
chain splitting.
Under Shapiro
the normalized
fiber averages
are precisely
the usual
restriction--corestriction
operators along
$H_9\subset H_3
\subset\Gamma$,
and the strict
nested identity
follows.

Finally the
marked index-nine
Schreier--Tietze
construction above
gives a genuine
minimal $P_9$
one-relator
Fox complex
after the
nine tame,
eight wild,
and eight
spanning-tree
cancellations.
The restricted
$H_9$ coefficient
action is pro-$2$
by hypothesis,
so this minimal
complex computes
$R\Gamma(H_9,B)$.
Its direct
summands then
give the
corresponding
derived objects,
and the constant
one computes
$R\Gamma(\Gamma,B)$
by split restriction
and the
$\frac19$
corestriction
projector.
The result
asserts a
selected
coefficient-sector
chain equivalence,
not a
presentation-independent
global group-algebra
syzygy frame.
\end{proof}

\begin{corollary}[Higher $3$-power tower projectors and faithful finite jets]
\label{cor:dyadic-f135-all-threepower-coherence}
For $n=3^k$ and
$0\le j\le k$,
let $P_{3^j|n}$
average over the
$n/3^j$ indices
congruent modulo
$3^j$:
\begin{equation}\label{eq:dyadic-f135-all-threepower-projectors}
(P_{3^j|n})_{ab}
=\begin{cases}
3^j/n,&a\equiv b\pmod{3^j},\\
0,&\text{otherwise}.
\end{cases}
\end{equation}
Every such matrix
has odd-unit
denominators,
is an integral
$\Gamma$-equivariant
idempotent,
and any two obey
\begin{equation}\label{eq:dyadic-f135-projector-semilattice}
\boxed{\qquad
P_{3^i|n}P_{3^j|n}
=P_{3^{\min(i,j)}|n}.
\qquad}
\end{equation}
The successive
differences of
these nested
idempotents are
pairwise orthogonal
and give a
strictly transitive
finite Hecke
decomposition of
the coinduced
$n$-sheet
Fox--Shapiro
complex.
On every chosen
finite pro-$2$
word and
coefficient jet,
the corresponding
marked Schreier
and Nielsen--Fox
comparison
terminates:
all $1/n$
powers are finite
integer powers
in finite
pro-$2$ quotients,
and all necessary
Fox pivots have
odd augmentation.
Thus this selected
$3$-primary odd-index
tower gives a
concrete family
of \emph{strict
integral descent
coherences},
not just a
single nine-sheet
numerical example.
\end{corollary}

\begin{proof}
The map from
$n$ cosets to
$3^j$ residue
classes is
$\Gamma$-equivariant.
Normalized fiber
sum is the
integral idempotent
$P_{3^j|n}$.
Composing any
two fiber
averages gives
the coarser
average, proving
\eqref{eq:dyadic-f135-projector-semilattice}.
Successive
differences are
therefore
orthogonal.
Their strict
cochain
commutation and
finite-jet
effectivity
follow from
the preceding
theorems and
the pro-$2$
unit Fox
minors.
\end{proof}

\begin{warning}[The precise scope of $3$-power tower closure]
\label{warn:dyadic-f135-threepower-scope}
The index-nine
example is
\emph{not} a
second copy of
the $S_3$
quotient.
Its tame
closure group
is
$C_9\rtimes C_6$,
and $\sigma_2$
acts by inversion,
while $\sigma^{-1}$
acts by
multiplication by
five on the
nine cosets.
Using the
three-coset
identity
$-i=\overline2i$
at nine cosets
would be incorrect
and would spoil
the marked group
relations.

The primitive
Schreier--Tietze
word transformations
are made in
chosen marked
free pro-$2$
coordinates.
Their inverse
limits are
effective at
each finite
quotient but
are not
presentation-independent
canonical
word formulas.
The theorem
supplies a
strict integral
odd-index
tower for
the specified
fields
$F_{3^k}$,
not a
functorial
minimal Fox
atlas for
all possible
non-pro-$2$
residual
quotients or
all local fields.
It also does
not assert
that the full
two-relator
group-algebra
complex is
an exact
universal
$\mathbf Z_2[[\Gamma]]$
resolution.
\end{warning}
